\section{AHSME 1979}
        \begin{enumerate}
        	\item (\textit{AHSME 1979}) 
                    \begin{figure}[H]
                        \centering
                        \includegraphics[width=0.5\textwidth]{Problem/AHSME_I_1979_1_Problem_Figure_1.png}
                    \end{figure}
                    


 If rectangle ABCD has area 72 square meters and E and G are the midpoints of sides AD and CD, respectively, then the area of rectangle DEFG in square meters is


 $\textbf{(A) }8\qquad \textbf{(B) }9\qquad \textbf{(C) }12\qquad \textbf{(D) }18\qquad \textbf{(E) }24$


	\item (\textit{AHSME 1979}) For all non-zero real numbers $x$ and $y$ such that $x-y=xy, \frac{1}{x}-\frac{1}{y}$ equals


 $\textbf{(A) }\frac{1}{xy}\qquad \textbf{(B) }\frac{1}{x-y}\qquad \textbf{(C) }0\qquad \textbf{(D) }-1\qquad \textbf{(E) }y-x$


	\item (\textit{AHSME 1979}) 
                    \begin{figure}[H]
                        \centering
                        \includegraphics[width=0.5\textwidth]{Problem/AHSME_I_1979_3_Problem_Figure_1.png}
                    \end{figure}
                    


 In the adjoining figure, $ABCD$ is a square, $ABE$ is an equilateral triangle and point $E$ is outside square $ABCD$. 
What is the measure of $\measuredangle AED$ in degrees?


 $\textbf{(A) }10\qquad \textbf{(B) }12.5\qquad \textbf{(C) }15\qquad \textbf{(D) }20\qquad \textbf{(E) }25$


	\item (\textit{AHSME 1979}) For all real numbers $x, x[x\{x(2-x)-4\}+10]+1=$
$\textbf{(A) }-x^4+2x^3+4x^2+10x+1\qquad \textbf{(B) }-x^4-2x^3+4x^2+10x+1\qquad\\ \textbf{(C) }-x^4-2x^3-4x^2+10x+1\qquad \textbf{(D) }-x^4-2x^3-4x^2-10x+1\qquad\\ \textbf{(E) }-x^4+2x^3-4x^2+10x+1$


	\item (\textit{AHSME 1979}) Find the sum of the digits of the largest even three digit number (in base ten representation) 
which is not changed when its units and hundreds digits are interchanged.


 $\textbf{(A) }22\qquad \textbf{(B) }23\qquad \textbf{(C) }24\qquad \textbf{(D) }25\qquad \textbf{(E) }26$


	\item (\textit{AHSME 1979}) $\frac{3}{2}+\frac{5}{4}+\frac{9}{8}+\frac{17}{16}+\frac{33}{32}+\frac{65}{64}-7=$
$\textbf{(A) }-\frac{1}{64}\qquad \textbf{(B) }-\frac{1}{16}\qquad \textbf{(C) }0\qquad \textbf{(D) }\frac{1}{16}\qquad \textbf{(E) }\frac{1}{64}$


	\item (\textit{AHSME 1979}) The square of an integer is called a perfect square. If $x$ is a perfect square, the next larger perfect square is


 $\textbf{(A) }x+1\qquad \textbf{(B) }x^2+1\qquad \textbf{(C) }x^2+2x+1\qquad \textbf{(D) }x^2+x\qquad \textbf{(E) }x+2\sqrt{x}+1$


	\item (\textit{AHSME 1979}) Find the area of the smallest region bounded by the graphs of $y=|x|$ and $x^2+y^2=4$.


 $\textbf{(A) }\frac{\pi}{4}\qquad \textbf{(B) }\frac{3\pi}{4}\qquad \textbf{(C) }\pi\qquad \textbf{(D) }\frac{3\pi}{2}\qquad \textbf{(E) }2\pi$


	\item (\textit{AHSME 1979}) The product of $\sqrt[3]{4}$ and $\sqrt[4]{8}$ equals


 $\textbf{(A) }\sqrt[7]{12}\qquad \textbf{(B) }2\sqrt[7]{12}\qquad \textbf{(C) }\sqrt[7]{32}\qquad \textbf{(D) }\sqrt[12]{32}\qquad \textbf{(E) }2\sqrt[12]{32}$


	\item (\textit{AHSME 1979}) If $P_1P_2P_3P_4P_5P_6$ is a regular hexagon whose apothem (distance from the center to midpoint of a side) is $2$, 
and $Q_i$ is the midpoint of side $P_iP_{i+1}$ for $i=1,2,3,4$, then the area of quadrilateral $Q_1Q_2Q_3Q_4$ is


 $\textbf{(A) }6\qquad \textbf{(B) }2\sqrt{6}\qquad \textbf{(C) }\frac{8\sqrt{3}}{3}\qquad \textbf{(D) }3\sqrt{3}\qquad \textbf{(E) }4\sqrt{3}$


	\item (\textit{AHSME 1979}) Find a positive integral solution to the equation
$\frac{1+3+5+\dots+(2n-1)}{2+4+6+\dots+2n}=\frac{115}{116}$
$\textbf{(A) }110\qquad \textbf{(B) }115\qquad \textbf{(C) }116\qquad \textbf{(D) }231\qquad\\ \textbf{(E) }\text{The equation has no positive integral solutions.}$


	\item (\textit{AHSME 1979}) 
                    \begin{figure}[H]
                        \centering
                        \includegraphics[width=0.5\textwidth]{Problem/AHSME_I_1979_12_Problem_Figure_1.png}
                    \end{figure}
                    


 In the adjoining figure, $CD$ is the diameter of a semi-circle with center $O$. Point $A$ lies on the extension of $DC$ past $C$; 
point $E$ lies on the semi-circle, and $B$ is the point of intersection (distinct from $E$ ) of line segment $AE$ with the semi-circle. 
If length $AB$ equals length $OD$, and the measure of $\measuredangle EOD$ is $45^\circ$, then the
measure of $\measuredangle BAO$ is


 $\textbf{(A) }10^\circ\qquad \textbf{(B) }15^\circ\qquad \textbf{(C) }20^\circ\qquad \textbf{(D) }25^\circ\qquad \textbf{(E) }30^\circ$


	\item (\textit{AHSME 1979}) The inequality $y-x<\sqrt{x^2}$ is satisfied if and only if


 $\textbf{(A) }y<0\text{ or }y<2x\text{ (or both inequalities hold)}\qquad \textbf{(B) }y>0\text{ or }y<2x\text{ (or both inequalities hold)}\qquad \textbf{(C) }y^2<2xy\qquad \textbf{(D) }y<0\qquad \textbf{(E) }x>0\text{ and }y<2x$


	\item (\textit{AHSME 1979}) In a certain sequence of numbers, the first number is $1$, and, for all $n\ge 2$, the product of the first $n$ numbers in the sequence is $n^2$. 
The sum of the third and the fifth numbers in the sequence is


 $\textbf{(A) }\frac{25}{9}\qquad \textbf{(B) }\frac{31}{15}\qquad \textbf{(C) }\frac{61}{16}\qquad \textbf{(D) }\frac{576}{225}\qquad \textbf{(E) }34$


	\item (\textit{AHSME 1979}) Two identical jars are filled with alcohol solutions, the ratio of the volume of alcohol to the volume of water being $p : 1$ 
in one jar and $q : 1$ in the other jar. If the entire contents of the two jars are mixed together, 
the ratio of the volume of alcohol to the volume of water in the mixture is


 $\textbf{(A) }\frac{p+q}{2}\qquad \textbf{(B) }\frac{p^2+q^2}{p+q}\qquad \textbf{(C) }\frac{2pq}{p+q}\qquad \textbf{(D) }\frac{2(p^2+pq+q^2)}{3(p+q)}\qquad \textbf{(E) }\frac{p+q+2pq}{p+q+2}$


	\item (\textit{AHSME 1979}) A circle with area $A_1$ is contained in the interior of a larger circle with area $A_1+A_2$. If the radius of the larger circle is $3$, 
and if $A_1 , A_2, A_1 + A_2$ is an arithmetic progression, then the radius of the smaller circle is


 $\textbf{(A) }\frac{\sqrt{3}}{2}\qquad \textbf{(B) }1\qquad \textbf{(C) }\frac{2}{\sqrt{3}}\qquad \textbf{(D) }\frac{3}{2}\qquad \textbf{(E) }\sqrt{3}$


	\item (\textit{AHSME 1979}) 
                    \begin{figure}[H]
                        \centering
                        \includegraphics[width=0.5\textwidth]{Problem/AHSME_I_1979_17_Problem_Figure_1.png}
                    \end{figure}
                    


 Points $A , B, C$, and $D$ are distinct and lie, in the given order, on a straight line. 
Line segments $AB, AC$, and $AD$ have lengths $x, y$, and $z$, respectively. 
If line segments $AB$ and $CD$ may be rotated about points $B$ and $C$, respectively, 
so that points $A$ and $D$ coincide, to form a triangle with positive area, 
then which of the following three inequalities must be satisfied?


 $\textbf{I. }x<\frac{z}{2}\qquad\\ \textbf{II. }y<x+\frac{z}{2}\qquad\\ \textbf{III. }y<\frac{z}{2}\qquad$
$\textbf{(A) }\textbf{I. }\text{only}\qquad \textbf{(B) }\textbf{II. }\text{only}\qquad \textbf{(C) }\textbf{I. }\text{and }\textbf{II. }\text{only}\qquad \textbf{(D) }\textbf{II. }\text{and }\textbf{III. }\text{only}\qquad \textbf{(E) }\textbf{I. },\textbf{II. },\text{and }\textbf{III. }$


	\item (\textit{AHSME 1979}) To the nearest thousandth, $\log_{10}2$ is $.301$ and $\log_{10}3$ is $.477$. 
Which of the following is the best approximation of $\log_5 10$?


 $\textbf{(A) }\frac{8}{7}\qquad \textbf{(B) }\frac{9}{7}\qquad \textbf{(C) }\frac{10}{7}\qquad \textbf{(D) }\frac{11}{7}\qquad \textbf{(E) }\frac{12}{7}$


	\item (\textit{AHSME 1979}) Find the sum of the squares of all real numbers satisfying the equation $x^{256}-256^{32}=0$.


 $\textbf{(A) }8\qquad \textbf{(B) }128\qquad \textbf{(C) }512\qquad \textbf{(D) }65,536\qquad \textbf{(E) }2(256^{32})$


	\item (\textit{AHSME 1979}) If $a=\frac{1}{2}$ and $(a+1)(b+1)=2$ then the radian measure of $\arctan a + \arctan b$ equals


 $\textbf{(A) }\frac{\pi}{2}\qquad \textbf{(B) }\frac{\pi}{3}\qquad \textbf{(C) }\frac{\pi}{4}\qquad \textbf{(D) }\frac{\pi}{5}\qquad \textbf{(E) }\frac{\pi}{6}$


	\item (\textit{AHSME 1979}) The length of the hypotenuse of a right triangle is $h$ , and the radius of the inscribed circle is $r$. 
The ratio of the area of the circle to the area of the triangle is


 $\textbf{(A) }\frac{\pi r}{h+2r}\qquad \textbf{(B) }\frac{\pi r}{h+r}\qquad \textbf{(C) }\frac{\pi}{2h+r}\qquad \textbf{(D) }\frac{\pi r^2}{r^2+h^2}\qquad \textbf{(E) }\text{none of these}$


	\item (\textit{AHSME 1979}) Find the number of pairs $(m, n)$ of integers which satisfy the equation $m^3 + 6m^2 + 5m = 27n^3 + 9n^2 + 9n + 1$.


 $\textbf{(A) }0\qquad \textbf{(B) }1\qquad \textbf{(C) }3\qquad \textbf{(D) }9\qquad \textbf{(E) }\infty$


	\item (\textit{AHSME 1979}) The edges of a regular tetrahedron with vertices $A ,~ B,~ C$, and $D$ each have length one. 
Find the least possible distance between a pair of points $P$ and $Q$, where $P$ is on edge $AB$ and $Q$ is on edge $CD$.


 


                    \begin{figure}[H]
                        \centering
                        \includegraphics[width=0.5\textwidth]{Problem/AHSME_I_1979_23_Problem_Figure_1.png}
                    \end{figure}
                    
$\textbf{(A) }\frac{1}{2}\qquad \textbf{(B) }\frac{3}{4}\qquad \textbf{(C) }\frac{\sqrt{2}}{2}\qquad \textbf{(D) }\frac{\sqrt{3}}{2}\qquad \textbf{(E) }\frac{\sqrt{3}}{3}$


	\item (\textit{AHSME 1979}) Sides $AB,~ BC$, and $CD$ of (simple*) quadrilateral $ABCD$ have lengths $4,~ 5$, and $20$, respectively. 
If vertex angles $B$ and $C$ are obtuse and $\sin C = - \cos B =\frac{3}{5}$, then side $AD$ has length


 
            \begin{itemize}
        \item A polygon is called “simple” if it is not self intersecting.
            \end{itemize}
            $\textbf{(A) }24\qquad \textbf{(B) }24.5\qquad \textbf{(C) }24.6\qquad \textbf{(D) }24.8\qquad \textbf{(E) }25$


	\item (\textit{AHSME 1979}) If $q_1 ( x )$ and $r_ 1$ are the quotient and remainder, respectively, when the polynomial $x^ 8$ is divided by 
$x + \frac{1}{2}$ , and if $q_ 2 ( x )$ and $r_2$ are the quotient and remainder, respectively, 
when $q_ 1 ( x )$ is divided by $x + \frac{1}{2}$, then find $r_2$.


 $\textbf{(A) }\frac{1}{256}\qquad \textbf{(B) }-\frac{1}{16}\qquad \textbf{(C) }1\qquad \textbf{(D) }-16\qquad \textbf{(E) }256$


	\item (\textit{AHSME 1979}) The function $f$ satisfies the functional equation $f(x) +f(y) = f(x + y ) - xy - 1$ for every pair $x,~ y$
of real numbers. If $f( 1) = 1$, then the number of integers $n \neq 1$ for which $f ( n ) = n$ is


 $\textbf{(A) }0\qquad \textbf{(B) }1\qquad \textbf{(C) }2\qquad \textbf{(D) }3\qquad \textbf{(E) }\infty$


	\item (\textit{AHSME 1979}) An ordered pair $( b , c )$ of integers, each of which has absolute value less than or equal to five, is chosen at random, with each
such ordered pair having an equal likelihood of being chosen. What is the probability that the equation $x^ 2 + bx + c = 0$ will
not have distinct positive real roots?


 $\textbf{(A) }\frac{106}{121}\qquad \textbf{(B) }\frac{108}{121}\qquad \textbf{(C) }\frac{110}{121}\qquad \textbf{(D) }\frac{112}{121}\qquad \textbf{(E) }\text{none of these}$


	\item (\textit{AHSME 1979}) 
                    \begin{figure}[H]
                        \centering
                        \includegraphics[width=0.5\textwidth]{Problem/AHSME_I_1979_28_Problem_Figure_1.png}
                    \end{figure}
                    


 Circles with centers $A ,~ B$, and $C$ each have radius $r$, where $1 < r < 2$. 
The distance between each pair of centers is $2$. If $B'$ is the point of intersection of circle $A$ and circle $C$ 
which is outside circle $B$, and if $C'$ is the point of intersection of circle $A$ and circle $B$ which is outside circle $C$, 
then length $B'C'$ equals


 $\textbf{(A) }3r-2\qquad \textbf{(B) }r^2\qquad \textbf{(C) }r+\sqrt{3(r-1)}\qquad\\ \textbf{(D) }1+\sqrt{3(r^2-1)}\qquad \textbf{(E) }\text{none of these}$


	\item (\textit{AHSME 1979}) For each positive number $x$, let 
$f(x)=\frac{\left(x+\frac{1}{x}\right)^6-\left(x^6+\frac{1}{x^6}\right)-2} {\left(x+\frac{1}{x}\right)^3+\left(x^3+\frac{1}{x^3}\right)}$. 
The minimum value of $f(x)$ is


 $\textbf{(A) }1\qquad \textbf{(B) }2\qquad \textbf{(C) }3\qquad \textbf{(D) }4\qquad \textbf{(E) }6$


	\item (\textit{AHSME 1979}) 
                    \begin{figure}[H]
                        \centering
                        \includegraphics[width=0.5\textwidth]{Problem/AHSME_I_1979_30_Problem_Figure_1.png}
                    \end{figure}
                    


 In $\triangle ABC$, $E$ is the midpoint of side $BC$ and $D$ is on side $AC$. 
If the length of $AC$ is $1$ and $\measuredangle BAC = 60^\circ, \measuredangle ABC = 100^\circ, \measuredangle ACB = 20^\circ$ and 
$\measuredangle DEC = 80^\circ$, then the area of $\triangle ABC$ plus twice the area of $\triangle CDE$ equals


 $\textbf{(A) }\frac{1}{4}\cos 10^\circ\qquad \textbf{(B) }\frac{\sqrt{3}}{8}\qquad \textbf{(C) }\frac{1}{4}\cos 40^\circ\qquad \textbf{(D) }\frac{1}{4}\cos 50^\circ\qquad \textbf{(E) }\frac{1}{8}$


\end{enumerate}